\documentclass[11pt,a4paper]{letter}
\usepackage[T1]{fontenc}
\usepackage{lmodern}
\usepackage[margin=2.5cm]{geometry}
\usepackage{xcolor}
\usepackage[hidelinks]{hyperref}
% Values awaiting the final runs or a co-author decision
\newcommand{\TBD}[1]{\textcolor{red}{\textbf{[#1]}}}

\signature{Keisuke Nishioka\\
Institute of Continuum Mechanics, Leibniz Universit\"at Hannover\\
\href{mailto:keisuke.nishioka@stud.uni-hannover.de}{keisuke.nishioka@stud.uni-hannover.de}\\
ORCID: \href{https://orcid.org/0009-0001-9360-1336}{0009-0001-9360-1336}}
\address{Institute of Continuum Mechanics\\Leibniz Universit\"at Hannover\\An der Universit\"at 1\\30823 Garbsen, Germany}
\date{\TBD{date}}

\begin{document}
\begin{letter}{Editors\\Bulletin of Mathematical Biology}

\opening{Dear Editors,}

On behalf of all co-authors, I submit the manuscript
``GPU-accelerated Bayesian inference of multi-species biofilm interaction parameters via TMCMC''
by K.~Nishioka, F.~Klempt, H.~Geisler, R.~Mukherjee, N.~Heine, K.~Doll-Nikutta, M.~Stiesch, S.~P.~Szafra\'{n}ski, M.~Soleimani and P.~Junker
for consideration as an Original Article in the \emph{Bulletin of Mathematical Biology}.

Peri-implantitis is driven by multi-species biofilms in which a commensal community shifts towards a dysbiotic one dominated by \emph{Porphyromonas gingivalis}.
Which interactions between the species sustain this shift is difficult to measure directly.
We infer the full $5\times5$ interaction matrix of a thermodynamically consistent continuum biofilm model, derived from an extended Hamilton principle, from in-vitro time courses of five oral taxa under four combinations of community state and cultivation (Heine et al., \emph{Front.\ Oral Health} 2025, several of whom are co-authors of this work).
No interaction is fixed a priori, and every inference stage is accepted only after explicit convergence checks across independent seeds.

The main biological result is a testable prediction.
The effective \emph{F.~nucleatum}--\emph{P.~gingivalis} coefficient is positive only under dysbiotic dynamic cultivation and indistinguishable from zero in both commensal conditions.
In posterior predictive simulations, removing \emph{F.~nucleatum} from the consortium suppresses the late \emph{P.~gingivalis} increase with probability \TBD{0.87--0.90}; this can be tested directly by a leave-one-species-out co-culture.

We believe the manuscript fits the scope of the \emph{Bulletin} because it connects a mechanistically derived mathematical model with quantitative data through Bayesian inference, and uses the inferred model to make a biological prediction rather than only to fit data.
Methodologically, we report how identifiability was assessed (seed agreement, prior dependence on a wide box, posterior mass at the bounds), and we provide the GPU implementation that makes the repeated runs affordable.

The manuscript has not been published and is not under consideration elsewhere.
All authors have approved the submission \TBD{confirm}.
Code, run configurations, convergence records and posterior samples are archived at Zenodo (\TBD{DOI}).
The authors declare no competing interests.

\TBD{optional: suggested reviewers, without conflicts of interest with any co-author}

\closing{Sincerely,}

\end{letter}
\end{document}
